\subsection{群胚的例子}

\begin{enumerate}
\def\labelenumi{\arabic{enumi})}
\item
  设 G 为群。以 \(\mathcal{G}\) 表示这样的箭图：其顶点集缩为单个元素，其箭集为 G。\(\mathcal{G}\) 赋有由 G 的合成律诱导的合成律后，是一个群胚，称为与群 G 伴随的群胚。
\item
  设 X 为集合。设 G 为箭图 \((\mathrm{X}, \mathrm{X} \times \mathrm{X}, \mathrm{pr}_{1}, \mathrm{pr}_{2})\) ；在 G 上定义合成律：令 \((x, x') \cdot (x', x'') = (x, x'')\) （若 \(x, x', x''\) 是 X 的元素）。这个律是结合的；对每个 \(x \in X\) ，\((x, x)\) 是 x 处的中性元；箭 \((x, x')\) 的逆是箭 \((x', x)\) 。这样定义的群胚称为集合 X 的偶对群胚。若 X 非空，它是单传递的。
\item
  设 X 与 S 为集合，\(p \colon \mathrm{X} \to \mathrm{S}\) 为映射。对 \(a \in \mathrm{S}\) ，记 \(\mathrm{X}_a = p^{-1}(a)\) 。对 S 中的 \(a\) 与 \(b\) ，设 \(\mathrm{G}_{a,b}\) 为集 \(\mathcal{B}(\mathrm{X}_a; \mathrm{X}_b)\) （从 \(\mathrm{X}_a\) 到 \(\mathrm{X}_b\) 的双射所成的），并设 G 为诸集 \(\mathrm{G}_{a,b}\) 的不交并。设 \(o\) 与 \(t\) 为从 G 到 S 的映射，使得 \(o(f) = a\) 与 \(t(f) = b\) （对每个元素 \(f \in \mathrm{G}_{a,b}\) ）。四元组 \((\mathrm{S}, \mathrm{G}, o, t)\) 是一个箭图。给它赋予由 \(m(f, g) = g \circ f\) （若 \(f \in \mathcal{B}(\mathrm{X}_a; \mathrm{X}_b)\) 且 \(g \in \mathcal{B}(\mathrm{X}_b; \mathrm{X}_c)\) ，其中 \(a, b, c\) 是 S 的元素）定义的合成律。这是一个群胚；记作 \(\mathcal{B}(\mathrm{X}, p)\) ，称为 X 相对于 \(p\) 的置换群胚。
\item
  设 \((\mathrm{G}_{i})_{i\in\mathrm{I}}\) 为群胚族。以 G 表示诸底箭图族的积箭图；对每个 \(i\in I\) ，设 \(pr_{i}\colon G\to G_{i}\) 为典范箭图态射。G 上存在唯一的合成律，使 G 成为群胚，且对每个 \(i \in I\) ，\(\mathrm{pr}_i\) 是群胚态射。设 \(f = (f_i)_{i \in I}\) 与 \(g = (g_i)_{i \in I}\) 为 G 的箭；它们可复合，当且仅当箭 \(f_i\) 与 \(g_i\) 可复合（对 \(i \in I\) ）。这时有 \(fg = (f_i g_i)_{i \in I}\) 。
\end{enumerate}

箭图 G 赋有这个合成律后，称为族 \((\mathrm{G}_{i})_{i\in\mathrm{I}}\) 的积群胚。它满足下列泛性质：对每个群胚 \(G'\) 与每个族 \((\varphi_{i})_{i\in\mathrm{I}}\) （其中 \(\varphi_{i}\) 是从 \(G'\) 到 \(G_{i}\) 的群胚态射），存在唯一的群胚态射 \(\varphi: G'\to G\) ，使得 \(\varphi_{i}=\operatorname{pr}_{i}\circ\varphi\) （对所有 \(i\in I\) ）。

\begin{enumerate}
\def\labelenumi{\arabic{enumi})}
\setcounter{enumi}{4}
\tightlist
\item
  设 \(((\mathrm{G}_{i})_{i\in\mathrm{I}},(\varphi_{i,j})_{i\prec j})\) 为群胚的归纳系，以右过滤预序集 \((\mathrm{I},\prec)\) 为指标集，其中 \(\varphi_{i,j}\) 是群胚态射。设 G 为这样的箭图：其顶点集为 \(\varinjlim_{i}\operatorname{Som}(\mathbf{G}_{i})\) ，箭集为 \(\varinjlim_{i}\operatorname{Fl}(\mathbf{G}_{i})\) ，源映射与靶映射分别为映射 \(\varinjlim_{i}o_{G_{i}}\) 与 \(\varinjlim_{i}t_{G_{i}}\) 。诸 \(G_{i}\) 的合成律在 G 上诱导一个合成律，使 G 成为群胚（参见 A, I, 第 114 页）。典范映射 \(\operatorname{Som}(\mathbf{G}_{i})\to\operatorname{Som}(\mathbf{G})\) 与 \(\operatorname{Fl}(\mathbf{G}_{i})\to\operatorname{Fl}(\mathbf{G})\) 定义了一个群胚态射 \(\varphi_{i}\) （从 \(G_{i}\) 到 G 的）。若 i、j 是 I 中满足 \(i\prec j\) 的元素，则有 \(\varphi_{j}\circ\varphi_{i,j}=\varphi_{i}\) 。群胚 G 称为群胚族 \(G_{i}\) 的归纳极限，记作 \(\varinjlim_{i}G_{i}\) 。它满足下列泛性质：对每个群胚 H 与每个族 \((\psi_{i})_{i\in\mathrm{I}}\) （其中对 \(i\in I\) ，\(\psi_{i}:G_{i}\to H\) 是群胚态射，并且 \(\psi_{j}\circ\varphi_{i,j}=\psi_{i}\) 对 I 的元素的每对 \((i,j)\) （满足 \(i\prec j\) ）成立），存在唯一的群胚态射 \(\psi:G\to H\) 使得 \(\psi_{i}=\psi\circ\varphi_{i}\) 。
\end{enumerate}

关于当不假设集 I 右过滤时满足这个泛性质的群胚的存在性，参见 II, 第 228 页, 习题 3。
% semantic-fragments: legacy-input-seam
\relax{}
